% ============================================================
%  ECON306: W10D2 — NZ Education System; Student Loan Scheme
%          Economics (Income-Contingent, Interest-Free)
%  Date: Thursday 24 September 2026
%  SOURCE: old w12_week12 (student loans section) from FULL deck
%    (lines 3700–3730) + Course Map W10D2 LOs.
% ============================================================

\documentclass[aspectratio=169, 12pt]{beamer}

% ============================================================
%  ECON306: Shared Preamble — The Economics of Health and Education
%  University of Otago — Semester 2, 2026
%
%  SHARED BY ALL ECON306 DAY DECKS.
%  Input this file immediately after \documentclass in each deck:
%    % ============================================================
%  ECON306: Shared Preamble — The Economics of Health and Education
%  University of Otago — Semester 2, 2026
%
%  SHARED BY ALL ECON306 DAY DECKS.
%  Input this file immediately after \documentclass in each deck:
%    % ============================================================
%  ECON306: Shared Preamble — The Economics of Health and Education
%  University of Otago — Semester 2, 2026
%
%  SHARED BY ALL ECON306 DAY DECKS.
%  Input this file immediately after \documentclass in each deck:
%    \input{../Preambles/econ306_preamble.tex}
%  Do NOT include \documentclass or per-deck metadata here.
% ============================================================

% ─── Theme & Colours ────────────────────────────────────────
\usetheme{default}
\usecolortheme{default}
\usefonttheme{professionalfonts}
\definecolor{OtagoBlue}{HTML}{003057}
\definecolor{OtagoGold}{HTML}{A8923A}
\definecolor{TextGrey}{HTML}{3A3A3A}
\definecolor{AccentBlue}{HTML}{2B6CB0}
\definecolor{LightBg}{HTML}{F7F9FC}

\setbeamercolor{frametitle}{fg=OtagoBlue, bg=white}
\setbeamercolor{title}{fg=OtagoBlue}
\setbeamercolor{subtitle}{fg=OtagoGold}
\setbeamercolor{author}{fg=TextGrey}
\setbeamercolor{date}{fg=TextGrey}
\setbeamercolor{institute}{fg=TextGrey}
\setbeamercolor{normal text}{fg=TextGrey, bg=white}
\setbeamercolor{item}{fg=OtagoBlue}
\setbeamercolor{subitem}{fg=AccentBlue}
\setbeamercolor{block title}{bg=OtagoBlue!12, fg=OtagoBlue}
\setbeamercolor{block body}{bg=OtagoBlue!5, fg=TextGrey}
\setbeamercolor{alertblock title}{bg=OtagoGold!20, fg=OtagoBlue}
\setbeamercolor{alertblock body}{bg=OtagoGold!8, fg=TextGrey}
\setbeamercolor{alerted text}{fg=OtagoGold!80!black}
\setbeamercolor{structure}{fg=OtagoBlue}

\setbeamertemplate{frametitle}{%
	\vspace{0.35em}%
	\textbf{\insertframetitle}\par%
	\vspace{0.05em}%
	\textcolor{OtagoGold}{\rule{\textwidth}{0.6pt}}}

\setbeamertemplate{navigation symbols}{}
\setbeamertemplate{footline}{%
	\leavevmode\hbox{%
		\begin{beamercolorbox}[wd=.40\paperwidth,ht=2.5ex,dp=1.5ex,leftskip=.3cm]{author in head/foot}%
			\textcolor{TextGrey}{\footnotesize ECON306 \textbullet\ University of Otago}%
		\end{beamercolorbox}%
		\begin{beamercolorbox}[wd=.47\paperwidth,ht=2.5ex,dp=1.5ex,center]{title in head/foot}%
			\textcolor{TextGrey}{\footnotesize \insertshorttitle}%
		\end{beamercolorbox}%
		\begin{beamercolorbox}[wd=.13\paperwidth,ht=2.5ex,dp=1.5ex,rightskip=.3cm,right]{date in head/foot}%
			\textcolor{TextGrey}{\footnotesize \insertframenumber\,/\,\inserttotalframenumber}%
	\end{beamercolorbox}}\vskip0pt}

\setbeamertemplate{itemize item}{\textcolor{OtagoGold}{\small$\blacktriangleright$}}
\setbeamertemplate{itemize subitem}{\textcolor{AccentBlue}{\small$\triangleright$}}
\setbeamertemplate{enumerate item}{\textcolor{OtagoBlue}{\small\bfseries\insertenumlabel.}}

\setbeamertemplate{title page}{%
	\begin{tikzpicture}[remember picture,overlay]
		\fill[OtagoBlue] (current page.north west)
		rectangle ([yshift=-0.38\paperheight]current page.north east);
	\end{tikzpicture}
	\vspace{0.06\paperheight}
	\begin{beamercolorbox}[wd=\textwidth, sep=0pt]{title}
		{\color{white}\Large\textbf{\inserttitle}}\\[0.4em]
		{\color{OtagoGold}\large\insertsubtitle}
	\end{beamercolorbox}
	\vspace{1.4em}
	\begin{beamercolorbox}[wd=\textwidth, sep=0pt]{author}
		{\color{TextGrey}\normalsize\insertauthor}\\[0.2em]
		{\color{TextGrey}\small\insertinstitute}\\[0.2em]
		{\color{TextGrey}\small\insertdate}
\end{beamercolorbox}}

% ─── Packages ──────────────────────────────────────────────
\usepackage{tikz}
\usepackage{booktabs}
\usepackage{array}
\usepackage{graphicx}
\usepackage{hyperref}
\usepackage{amsmath}
\usepackage{amssymb}
\usepackage{colortbl}
\usepackage{multirow}
\usepackage{xcolor}
\hypersetup{colorlinks=true, linkcolor=AccentBlue, urlcolor=AccentBlue,
	citecolor=AccentBlue, pdfauthor={Austin Henderson},
	pdftitle={ECON306 S2 2026}}

% ─── Images folder ─────────────────────────────────────────
%  Place all slide images in the 'images/' subfolder next to each deck.
%  Usage: \includegraphics[width=0.85\textwidth]{figure_name}
\graphicspath{{images/}}

% ─── Reusable macros ───────────────────────────────────────
% Recap slide at start of each lecture
\newcommand{\recapslide}[1]{%
	\begin{frame}{Recap: What Did We Cover Last Time?}
		\begin{block}{Key points from last lecture}
			\begin{itemize}#1\end{itemize}
		\end{block}
\end{frame}}

% Plan slide at start of each lecture
\newcommand{\planslide}[2]{%
	\begin{frame}{Today's Plan — #1}
		\begin{columns}
			\column{0.6\textwidth}
			\begin{itemize}#2\end{itemize}
			\column{0.38\textwidth}
			\begin{alertblock}{Reading}
				\footnotesize See Brightspace (Aoroa) for required readings before this lecture.
			\end{alertblock}
		\end{columns}
\end{frame}}

% Summary slide at end of each lecture
\newcommand{\summaryside}[1]{%
	\begin{frame}{Lecture Summary}
		\begin{exampleblock}{Key takeaways}
			\begin{itemize}#1\end{itemize}
		\end{exampleblock}
		\vspace{0.5em}
		{\footnotesize\textcolor{OtagoGold}{\textbf{Brightspace (Aoroa):}} slides, readings, and any announcements posted there.}
\end{frame}}

% NZ spotlight box
\newcommand{\nzbox}[2]{%
	\begin{alertblock}{\includegraphics[height=0.8em]{nz_flag_icon}\hspace{0.3em} NZ Spotlight: #1}
		#2
\end{alertblock}}

% Tutorial callout
\newcommand{\tutorialbox}[1]{%
	\begin{block}{\textbf{Tutorial This Week}}
		#1
\end{block}}

% ─── Section dividers ──────────────────────────────────────
\AtBeginSection[]{
	\begin{frame}[plain]
		\begin{tikzpicture}[remember picture,overlay]
			\fill[OtagoBlue!95] (current page.north west)
			rectangle (current page.south east);
			\node[anchor=center, text=white, font=\Large\bfseries,
			text width=0.9\paperwidth, align=center]
			at (current page.center) {\insertsectionhead};
			\node[anchor=south, text=OtagoGold, font=\normalsize,
			text width=0.9\paperwidth, align=center]
			at ([yshift=2.5cm]current page.south) {ECON306 — University of Otago};
		\end{tikzpicture}
\end{frame}}

%  Do NOT include \documentclass or per-deck metadata here.
% ============================================================

% ─── Theme & Colours ────────────────────────────────────────
\usetheme{default}
\usecolortheme{default}
\usefonttheme{professionalfonts}
\definecolor{OtagoBlue}{HTML}{003057}
\definecolor{OtagoGold}{HTML}{A8923A}
\definecolor{TextGrey}{HTML}{3A3A3A}
\definecolor{AccentBlue}{HTML}{2B6CB0}
\definecolor{LightBg}{HTML}{F7F9FC}

\setbeamercolor{frametitle}{fg=OtagoBlue, bg=white}
\setbeamercolor{title}{fg=OtagoBlue}
\setbeamercolor{subtitle}{fg=OtagoGold}
\setbeamercolor{author}{fg=TextGrey}
\setbeamercolor{date}{fg=TextGrey}
\setbeamercolor{institute}{fg=TextGrey}
\setbeamercolor{normal text}{fg=TextGrey, bg=white}
\setbeamercolor{item}{fg=OtagoBlue}
\setbeamercolor{subitem}{fg=AccentBlue}
\setbeamercolor{block title}{bg=OtagoBlue!12, fg=OtagoBlue}
\setbeamercolor{block body}{bg=OtagoBlue!5, fg=TextGrey}
\setbeamercolor{alertblock title}{bg=OtagoGold!20, fg=OtagoBlue}
\setbeamercolor{alertblock body}{bg=OtagoGold!8, fg=TextGrey}
\setbeamercolor{alerted text}{fg=OtagoGold!80!black}
\setbeamercolor{structure}{fg=OtagoBlue}

\setbeamertemplate{frametitle}{%
	\vspace{0.35em}%
	\textbf{\insertframetitle}\par%
	\vspace{0.05em}%
	\textcolor{OtagoGold}{\rule{\textwidth}{0.6pt}}}

\setbeamertemplate{navigation symbols}{}
\setbeamertemplate{footline}{%
	\leavevmode\hbox{%
		\begin{beamercolorbox}[wd=.40\paperwidth,ht=2.5ex,dp=1.5ex,leftskip=.3cm]{author in head/foot}%
			\textcolor{TextGrey}{\footnotesize ECON306 \textbullet\ University of Otago}%
		\end{beamercolorbox}%
		\begin{beamercolorbox}[wd=.47\paperwidth,ht=2.5ex,dp=1.5ex,center]{title in head/foot}%
			\textcolor{TextGrey}{\footnotesize \insertshorttitle}%
		\end{beamercolorbox}%
		\begin{beamercolorbox}[wd=.13\paperwidth,ht=2.5ex,dp=1.5ex,rightskip=.3cm,right]{date in head/foot}%
			\textcolor{TextGrey}{\footnotesize \insertframenumber\,/\,\inserttotalframenumber}%
	\end{beamercolorbox}}\vskip0pt}

\setbeamertemplate{itemize item}{\textcolor{OtagoGold}{\small$\blacktriangleright$}}
\setbeamertemplate{itemize subitem}{\textcolor{AccentBlue}{\small$\triangleright$}}
\setbeamertemplate{enumerate item}{\textcolor{OtagoBlue}{\small\bfseries\insertenumlabel.}}

\setbeamertemplate{title page}{%
	\begin{tikzpicture}[remember picture,overlay]
		\fill[OtagoBlue] (current page.north west)
		rectangle ([yshift=-0.38\paperheight]current page.north east);
	\end{tikzpicture}
	\vspace{0.06\paperheight}
	\begin{beamercolorbox}[wd=\textwidth, sep=0pt]{title}
		{\color{white}\Large\textbf{\inserttitle}}\\[0.4em]
		{\color{OtagoGold}\large\insertsubtitle}
	\end{beamercolorbox}
	\vspace{1.4em}
	\begin{beamercolorbox}[wd=\textwidth, sep=0pt]{author}
		{\color{TextGrey}\normalsize\insertauthor}\\[0.2em]
		{\color{TextGrey}\small\insertinstitute}\\[0.2em]
		{\color{TextGrey}\small\insertdate}
\end{beamercolorbox}}

% ─── Packages ──────────────────────────────────────────────
\usepackage{tikz}
\usepackage{booktabs}
\usepackage{array}
\usepackage{graphicx}
\usepackage{hyperref}
\usepackage{amsmath}
\usepackage{amssymb}
\usepackage{colortbl}
\usepackage{multirow}
\usepackage{xcolor}
\hypersetup{colorlinks=true, linkcolor=AccentBlue, urlcolor=AccentBlue,
	citecolor=AccentBlue, pdfauthor={Austin Henderson},
	pdftitle={ECON306 S2 2026}}

% ─── Images folder ─────────────────────────────────────────
%  Place all slide images in the 'images/' subfolder next to each deck.
%  Usage: \includegraphics[width=0.85\textwidth]{figure_name}
\graphicspath{{images/}}

% ─── Reusable macros ───────────────────────────────────────
% Recap slide at start of each lecture
\newcommand{\recapslide}[1]{%
	\begin{frame}{Recap: What Did We Cover Last Time?}
		\begin{block}{Key points from last lecture}
			\begin{itemize}#1\end{itemize}
		\end{block}
\end{frame}}

% Plan slide at start of each lecture
\newcommand{\planslide}[2]{%
	\begin{frame}{Today's Plan — #1}
		\begin{columns}
			\column{0.6\textwidth}
			\begin{itemize}#2\end{itemize}
			\column{0.38\textwidth}
			\begin{alertblock}{Reading}
				\footnotesize See Brightspace (Aoroa) for required readings before this lecture.
			\end{alertblock}
		\end{columns}
\end{frame}}

% Summary slide at end of each lecture
\newcommand{\summaryside}[1]{%
	\begin{frame}{Lecture Summary}
		\begin{exampleblock}{Key takeaways}
			\begin{itemize}#1\end{itemize}
		\end{exampleblock}
		\vspace{0.5em}
		{\footnotesize\textcolor{OtagoGold}{\textbf{Brightspace (Aoroa):}} slides, readings, and any announcements posted there.}
\end{frame}}

% NZ spotlight box
\newcommand{\nzbox}[2]{%
	\begin{alertblock}{\includegraphics[height=0.8em]{nz_flag_icon}\hspace{0.3em} NZ Spotlight: #1}
		#2
\end{alertblock}}

% Tutorial callout
\newcommand{\tutorialbox}[1]{%
	\begin{block}{\textbf{Tutorial This Week}}
		#1
\end{block}}

% ─── Section dividers ──────────────────────────────────────
\AtBeginSection[]{
	\begin{frame}[plain]
		\begin{tikzpicture}[remember picture,overlay]
			\fill[OtagoBlue!95] (current page.north west)
			rectangle (current page.south east);
			\node[anchor=center, text=white, font=\Large\bfseries,
			text width=0.9\paperwidth, align=center]
			at (current page.center) {\insertsectionhead};
			\node[anchor=south, text=OtagoGold, font=\normalsize,
			text width=0.9\paperwidth, align=center]
			at ([yshift=2.5cm]current page.south) {ECON306 — University of Otago};
		\end{tikzpicture}
\end{frame}}

%  Do NOT include \documentclass or per-deck metadata here.
% ============================================================

% ─── Theme & Colours ────────────────────────────────────────
\usetheme{default}
\usecolortheme{default}
\usefonttheme{professionalfonts}
\definecolor{OtagoBlue}{HTML}{003057}
\definecolor{OtagoGold}{HTML}{A8923A}
\definecolor{TextGrey}{HTML}{3A3A3A}
\definecolor{AccentBlue}{HTML}{2B6CB0}
\definecolor{LightBg}{HTML}{F7F9FC}

\setbeamercolor{frametitle}{fg=OtagoBlue, bg=white}
\setbeamercolor{title}{fg=OtagoBlue}
\setbeamercolor{subtitle}{fg=OtagoGold}
\setbeamercolor{author}{fg=TextGrey}
\setbeamercolor{date}{fg=TextGrey}
\setbeamercolor{institute}{fg=TextGrey}
\setbeamercolor{normal text}{fg=TextGrey, bg=white}
\setbeamercolor{item}{fg=OtagoBlue}
\setbeamercolor{subitem}{fg=AccentBlue}
\setbeamercolor{block title}{bg=OtagoBlue!12, fg=OtagoBlue}
\setbeamercolor{block body}{bg=OtagoBlue!5, fg=TextGrey}
\setbeamercolor{alertblock title}{bg=OtagoGold!20, fg=OtagoBlue}
\setbeamercolor{alertblock body}{bg=OtagoGold!8, fg=TextGrey}
\setbeamercolor{alerted text}{fg=OtagoGold!80!black}
\setbeamercolor{structure}{fg=OtagoBlue}

\setbeamertemplate{frametitle}{%
	\vspace{0.35em}%
	\textbf{\insertframetitle}\par%
	\vspace{0.05em}%
	\textcolor{OtagoGold}{\rule{\textwidth}{0.6pt}}}

\setbeamertemplate{navigation symbols}{}
\setbeamertemplate{footline}{%
	\leavevmode\hbox{%
		\begin{beamercolorbox}[wd=.40\paperwidth,ht=2.5ex,dp=1.5ex,leftskip=.3cm]{author in head/foot}%
			\textcolor{TextGrey}{\footnotesize ECON306 \textbullet\ University of Otago}%
		\end{beamercolorbox}%
		\begin{beamercolorbox}[wd=.47\paperwidth,ht=2.5ex,dp=1.5ex,center]{title in head/foot}%
			\textcolor{TextGrey}{\footnotesize \insertshorttitle}%
		\end{beamercolorbox}%
		\begin{beamercolorbox}[wd=.13\paperwidth,ht=2.5ex,dp=1.5ex,rightskip=.3cm,right]{date in head/foot}%
			\textcolor{TextGrey}{\footnotesize \insertframenumber\,/\,\inserttotalframenumber}%
	\end{beamercolorbox}}\vskip0pt}

\setbeamertemplate{itemize item}{\textcolor{OtagoGold}{\small$\blacktriangleright$}}
\setbeamertemplate{itemize subitem}{\textcolor{AccentBlue}{\small$\triangleright$}}
\setbeamertemplate{enumerate item}{\textcolor{OtagoBlue}{\small\bfseries\insertenumlabel.}}

\setbeamertemplate{title page}{%
	\begin{tikzpicture}[remember picture,overlay]
		\fill[OtagoBlue] (current page.north west)
		rectangle ([yshift=-0.38\paperheight]current page.north east);
	\end{tikzpicture}
	\vspace{0.06\paperheight}
	\begin{beamercolorbox}[wd=\textwidth, sep=0pt]{title}
		{\color{white}\Large\textbf{\inserttitle}}\\[0.4em]
		{\color{OtagoGold}\large\insertsubtitle}
	\end{beamercolorbox}
	\vspace{1.4em}
	\begin{beamercolorbox}[wd=\textwidth, sep=0pt]{author}
		{\color{TextGrey}\normalsize\insertauthor}\\[0.2em]
		{\color{TextGrey}\small\insertinstitute}\\[0.2em]
		{\color{TextGrey}\small\insertdate}
\end{beamercolorbox}}

% ─── Packages ──────────────────────────────────────────────
\usepackage{tikz}
\usepackage{booktabs}
\usepackage{array}
\usepackage{graphicx}
\usepackage{hyperref}
\usepackage{amsmath}
\usepackage{amssymb}
\usepackage{colortbl}
\usepackage{multirow}
\usepackage{xcolor}
\hypersetup{colorlinks=true, linkcolor=AccentBlue, urlcolor=AccentBlue,
	citecolor=AccentBlue, pdfauthor={Austin Henderson},
	pdftitle={ECON306 S2 2026}}

% ─── Images folder ─────────────────────────────────────────
%  Place all slide images in the 'images/' subfolder next to each deck.
%  Usage: \includegraphics[width=0.85\textwidth]{figure_name}
\graphicspath{{images/}}

% ─── Reusable macros ───────────────────────────────────────
% Recap slide at start of each lecture
\newcommand{\recapslide}[1]{%
	\begin{frame}{Recap: What Did We Cover Last Time?}
		\begin{block}{Key points from last lecture}
			\begin{itemize}#1\end{itemize}
		\end{block}
\end{frame}}

% Plan slide at start of each lecture
\newcommand{\planslide}[2]{%
	\begin{frame}{Today's Plan — #1}
		\begin{columns}
			\column{0.6\textwidth}
			\begin{itemize}#2\end{itemize}
			\column{0.38\textwidth}
			\begin{alertblock}{Reading}
				\footnotesize See Brightspace (Aoroa) for required readings before this lecture.
			\end{alertblock}
		\end{columns}
\end{frame}}

% Summary slide at end of each lecture
\newcommand{\summaryside}[1]{%
	\begin{frame}{Lecture Summary}
		\begin{exampleblock}{Key takeaways}
			\begin{itemize}#1\end{itemize}
		\end{exampleblock}
		\vspace{0.5em}
		{\footnotesize\textcolor{OtagoGold}{\textbf{Brightspace (Aoroa):}} slides, readings, and any announcements posted there.}
\end{frame}}

% NZ spotlight box
\newcommand{\nzbox}[2]{%
	\begin{alertblock}{\includegraphics[height=0.8em]{nz_flag_icon}\hspace{0.3em} NZ Spotlight: #1}
		#2
\end{alertblock}}

% Tutorial callout
\newcommand{\tutorialbox}[1]{%
	\begin{block}{\textbf{Tutorial This Week}}
		#1
\end{block}}

% ─── Section dividers ──────────────────────────────────────
\AtBeginSection[]{
	\begin{frame}[plain]
		\begin{tikzpicture}[remember picture,overlay]
			\fill[OtagoBlue!95] (current page.north west)
			rectangle (current page.south east);
			\node[anchor=center, text=white, font=\Large\bfseries,
			text width=0.9\paperwidth, align=center]
			at (current page.center) {\insertsectionhead};
			\node[anchor=south, text=OtagoGold, font=\normalsize,
			text width=0.9\paperwidth, align=center]
			at ([yshift=2.5cm]current page.south) {ECON306 — University of Otago};
		\end{tikzpicture}
\end{frame}}


\title[ECON306: W10D2 — NZ Education and Student Loans]{ECON306\\The Economics of Health and Education}
\subtitle{W10D2: NZ Education System; Student Loan Scheme Economics}
\author{Austin Henderson}
\institute{Department of Economics\\University of Otago}
\date{Thursday 24 September 2026}

\begin{document}

\begin{frame}[plain]
  \titlepage
\end{frame}

\recapslide{
  \item Education market failures: externalities, capital market failure, imperfect information
  \item Six characteristics framework maps to education as well as health
  \item NZ: free compulsory schooling Years 1--10; subsidised tertiary; income-contingent student loans
}

\planslide{NZ Student Loan Scheme: Economics and Equity}{
  \item The economics of income-contingent lending: solving capital market failure?
  \item NZ Student Loan Scheme: design, scale, and costs
  \item Interest-free lending: efficiency and equity implications
  \item The ``free first year'' policy: economic assessment
  \item Lifelong learning and whether the current scheme supports older workers
}

\begin{frame}{The Capital Market Failure in Education}
  \begin{columns}
    \column{0.52\textwidth}
      \footnotesize \textbf{Why can't students borrow freely?}
      \begin{itemize}
        \item Human capital cannot be collateralised (you cannot mortgage your degree)
        \item Banks reluctant to lend for unsecured investments with uncertain returns
        \item \textbf{Under-investment} results, especially for low-income families
        \item This is a \textbf{pure capital market failure}, not a subsidy argument
      \end{itemize}
    \column{0.45\textwidth}
      \begin{alertblock}{Policy response: income-contingent loans}
        \footnotesize
        If government provides loans repayable as a \textbf{percentage of future income},
        it removes repayment risk during low-income periods.\\[0.3em]
        This solves the collateral problem: repayment scales with the realised benefit.
      \end{alertblock}
  \end{columns}
\end{frame}

\begin{frame}{NZ Student Loan Scheme: Design and Scale}
  \begin{columns}
    \column{0.55\textwidth}
      \footnotesize
      \begin{itemize}
        \item Introduced 1992 to replace grants with loans; designed as income-contingent
        \item Total outstanding debt: $\approx$NZ\$9.8 billion (2024); $\approx$700,000 borrowers
        \item Average debt $\approx$\$22,000 at graduation
        \item \textbf{Income-contingent repayment}: 12\% of income above threshold
              ($\approx$\$22,828/year); zero repayments while below threshold
        \item \textbf{Interest-free} for NZ residents (introduced 2006)
        \item Estimated \textbf{subsidy cost}: $\approx$\$250M/year (Treasury)
      \end{itemize}
    \column{0.42\textwidth}
      \footnotesize \textbf{Key economic questions}
      \begin{itemize}
        \item Does income-contingency solve the capital market failure?
        \item Is interest-free lending efficient?
        \item What is the private return to tertiary education in NZ?
        \item Does the scheme achieve equity goals?
      \end{itemize}
  \end{columns}
\end{frame}

\begin{frame}{Is Interest-Free Lending Efficient?}
  \begin{columns}
    \column{0.52\textwidth}
      \footnotesize \textbf{Arguments for interest-free}
      \begin{itemize}
        \item Reduces disincentive to borrow among low-income students
        \item Corrects for \emph{additional} risk aversion in the borrowing decision
        \item Promotes equity: less deterrent for disadvantaged students
        \item Investment in human capital benefits society (externalities)
      \end{itemize}
    \column{0.45\textwidth}
      \begin{alertblock}{Arguments against (efficiency)}
        \footnotesize
        \begin{itemize}
          \item Capital market failure already solved by income-contingency; interest-free
                is an \emph{additional} subsidy on top
          \item Subsidy ($\approx$\$250M/year) disproportionately benefits graduates
                (above-average earners) --- regressive long-run
          \item May inflate demand for tertiary education beyond socially optimal
          \item Encourages borrowing for living costs, not just tuition
        \end{itemize}
      \end{alertblock}
  \end{columns}
\end{frame}

\begin{frame}{The ``Free First Year'' Policy}
  \small
  \begin{itemize}
    \item Introduced in NZ in 2018: first year of tertiary education fully funded by
          government for domestic students
    \item \textbf{Rationale}: lower the entry barrier; encourage more NZers to undertake
          post-secondary study; promote equity
    \item \textbf{Economic critique}:
    \begin{itemize}
      \item The binding constraint for most students was not fees for the first year ---
            it was uncertainty about completing the degree and returns thereafter
      \item Benefits accrue to students who would have enrolled anyway (deadweight loss)
      \item Significant fiscal cost for limited marginal impact on enrolments
      \item Does not address the equity gap for students with no family support for
            living costs in the first year
    \end{itemize}
  \end{itemize}
\end{frame}

\begin{frame}{The ``Free First Year'' Policy: NZ Evidence}
  \begin{exampleblock}{NZ evidence}
    Post-2018 enrolment data show modest increases, concentrated among higher-income groups
    rather than the disadvantaged students the policy targeted.
  \end{exampleblock}
  \vspace{0.5em}
  \small
  \textbf{Verdict}: The policy was a broad subsidy that largely benefited those who would
  have enrolled regardless. Better-targeted interventions (living cost support, mentoring
  for first-generation students) would deliver greater equity gains per dollar spent.
\end{frame}

\begin{frame}{Student Loans and Lifelong Learning}
  \footnotesize
  \begin{itemize}
    \item Traditional model: \textbf{education early in life $\Rightarrow$ work $\Rightarrow$ retirement}
    \item This model is becoming obsolete as populations age and labour markets change rapidly
    \item \textbf{NZ student loan scheme} was designed for young, full-time students:
    \begin{itemize}
      \item Living costs support assumes you are studying full-time
      \item Repayment schedule assumes long future working life to amortise debt
      \item Older workers upskilling mid-career have shorter repayment windows
    \end{itemize}
    \item \textbf{NZ's Tertiary Education Commission} funds vocational training and industry
          partnerships, but uptake by 45+ workers is limited
    \item \textbf{Human capital model prediction} (Week 10D3): returns to education fall with age
          (fewer working years to amortise investment) --- but longer working lives due to
          population ageing partially offset this
  \end{itemize}
  \begin{alertblock}{Policy gap}
    \footnotesize NZ lacks a well-designed lifelong learning financing system for mid-career workers.
  \end{alertblock}
\end{frame}

\begin{frame}{Equity in Tertiary Education Funding}
  \begin{columns}
    \column{0.52\textwidth}
      \footnotesize
      \begin{itemize}
        \item \textbf{User-pays case}: private returns high ($\approx$60--80\% wage premium);
              those who benefit should contribute
        \item \textbf{Subsidy case}: positive externalities mean private demand below socially optimal
        \item \textbf{Income-contingent loans}: solve capital market failure;
              risk-sharing between student and state
        \item \textbf{M\=aori and Pacific students}: additional barriers --- fewer family resources;
              returns underestimated due to labour market discrimination
      \end{itemize}
    \column{0.44\textwidth}
      \begin{alertblock}{Tutorial 6 preview (Week 12)}
        \footnotesize
        Applying MCDM to your own life: what do you value, what to do, where to live (1000Minds).\\[0.2em]
        Mon 5 Oct / Fri 9 Oct 2026.\\[0.2em]
        \textbf{Assignment due Mon 5 Oct, 9:00 am}.
      \end{alertblock}
  \end{columns}
\end{frame}

\summaryside{
  \footnotesize
  \item Capital market failure: students cannot collateralise human capital; income-contingent loans solve this
  \item NZ Student Loan Scheme: $\approx$\$9.8B outstanding; interest-free ($\approx$\$250M/year subsidy)
  \item Interest-free lending adds a subsidy beyond capital market correction; distributional debate (regressive long-run?)
  \item ``Free first year'' had limited equity impact; NZ lacks well-designed mid-career lifelong learning financing
}

\end{document}
